% Research note for the parent agent. This is a section fragment, not a standalone document.
\section{Predecessor from equational induction alone}
\label{sec:eq-predecessor}

This section removes a possible circularity from the restriction of existential
witnesses. The predecessor term is constructed before any general term-naming
principle, and without numerical coverage, a separation axiom, existential
introduction, or a case split under a hypothesis such as $c=0$.

\paragraph{The equational interface.}
Use the term constructors $0$, $S$, and
$I(n,a;u.h)$, the two computation equations
\[
 I(0,a;u.h)=a,\qquad
 I(Sn,a;u.h)=h[I(n,a;u.h)/u],
\]
ordinary equational reasoning, and induction for a literal equation. Thus,
if $z$ is absent from the parameters and hypotheses $\Gamma$, then
\[
 \frac{\Gamma\vdash E(0)\qquad
       \Gamma,E(z)\vdash E(Sz)}{\Gamma\vdash E(t)}.
\]
The unconditional cases rule
\[
 \frac{\Gamma\vdash E(0)\qquad\Gamma\vdash E(Sz)}
      {\Gamma\vdash E(t)}
\]
is the special case that ignores the induction hypothesis. This rule does
\emph{not} replace or specialize hypotheses of $\Gamma$ containing an
arbitrary term $c$. Every use of cases below proves an unconditional identity
with the displayed case variable and the other parameters independent.

All definitions below are abbreviations for scalar iterator terms. Write
\[
 C(c,a,b)=I(c,a;u.b),\qquad
 x+y=I(y,x;u.Su),\qquad 1=S0.
\]
Then $C(0,a,b)=a$, $C(Sz,a,b)=b$,
$x+0=x$, and $x+Sy=S(x+y)$. Equation induction gives
$0+x=x$, $Sx+y=S(x+y)$, associativity and commutativity of addition.
For example the second identity is induction on $y$, commutativity is
induction on $y$ using that identity, and associativity is induction on its
last argument. Hence finite sums can be rearranged, and $Sx=x+1$.
The notation $2x$ and $3x$ means $x+x$ and $x+x+x$.

\paragraph{Zero test and parity.}
Put
\[
 Z(x)=C(x,1,0),\qquad P(n)=I(n,0;u.Z(u)).
\]
The equations $Z(0)=1$, $Z(Sx)=0$, and
\[
 Z(Z(Zx))=Zx
\]
follow by computation and unconditional cases on $x$. Consequently
\[
 P(0)=0,\quad P(Sn)=Z(Pn),\quad Z(Z(Pn))=Pn.
\]
The last equation is induction on $n$: the base computes, and the step
is the displayed triple-zero-test identity. It follows that
$P(n+2)=Pn$. Induction on $n$ now gives
\[
 P(2n)=0,\qquad P(3n)=Pn.
\]
For the first step use $2(Sn)=2n+2$; for the second use
$3(Sn)=3n+3$ and $P(m+3)=Z(Pm)$.

\paragraph{Boolean interpolation without conditional cases.}
For any term template $H(v)$, with arbitrary other parameters, unconditional
cases on $y$ prove
\begin{equation}
 H(Z(Zy))=C(y,H(0),H(1)).
 \tag{BI}\label{eq:boolean-interpolation}
\end{equation}
Indeed $Z(Z0)=0$ and $Z(Z(Sz))=1$. Thus if an equation $p=Z(Zp)$ is
already available, equational congruence turns \eqref{eq:boolean-interpolation}
into $H(p)=C(p,H(0),H(1))$. In particular an identity between two term
templates at $0$ and at $1$ holds at such a $p$. This uses a previously proved
unconditional identity; it never substitutes into a contextual hypothesis.

\paragraph{The three-cycle and the ramp.}
Define
\[
 c(u)=Z(u)+2P(u),\qquad M(n)=I(n,0;u.c(u)),\qquad
 J(k,a)=I(k,a;u.c(u)).
\]
Computation gives $c(0)=1$, $c(1)=2$, and $c(2)=0$.
The iterator equations give $M(0)=0$, $M(Sn)=c(Mn)$, and induction on $k$
gives the addition law
\begin{equation}
 M(r+k)=J(k,M(r)).\tag{M+}\label{eq:mod3-add}
\end{equation}
Its step is $M(r+Sk)=c(M(r+k))=c(J(k,Mr))=J(Sk,Mr)$.

Now put
\[
 R(n)=I(n,0;u.S(u+M(u))).
\]
We claim $M(Rn)=Pn$. For its induction step write $r=Rn$ and $p=Pn$,
and use the induction hypothesis $Mr=p$. By \eqref{eq:mod3-add},
\[
 M(R(Sn))=c(M(r+Mr))=c(J(Mr,Mr))=c(J(p,p)).
\]
For a Boolean $p$, the identity
\[
 c(J(p,p))=Zp
\]
follows from (BI): at $p=0$ both sides are $1$, and at $p=1$ both sides
are $0$, since $c(J(1,1))=c(c(1))=c(2)=0$.
The already proved equation $p=Z(Zp)$ applies. This proves the induction
step because $P(Sn)=Zp$. The base is immediate.
We therefore have the simpler ramp recurrence
\begin{equation}
 R(0)=0,\qquad R(Sn)=Rn+Pn+1.\tag{R}\label{eq:ramp-rec}
\end{equation}

Induction now gives
\begin{equation}
 R(2n)=3n,\qquad R(3n)=3Rn+Pn.\tag{R23}\label{eq:ramp-multiples}
\end{equation}
For the first step, $P(2n)=0$ and $P(2n+1)=1$, so two applications of
\eqref{eq:ramp-rec} increase $R(2n)$ by $3$.
For the second, put $r=Rn$, $p=Pn$, and $q=Zp$.
The three successive parities at $3n$, $3n+1$, $3n+2$ are $p,q,p$.
Using the induction hypothesis, the next left side is therefore
\[
 R(3n+3)=3r+p+3+p+q+p=3r+3p+q+3,
\]
which is $3R(Sn)+P(Sn)=3(r+p+1)+q$. No conditional parity split is used.

\paragraph{A factor map.}
Let $H(x)=C(Px,Rx,x)$. Then
\[
 H(2x)=3x,\qquad H(1)=1,\qquad H(3x)=3H(x).
\]
The first two follow immediately. For the third put $p=Px$, $r=Rx$ and
use \eqref{eq:ramp-multiples} to obtain
\[
 H(3x)=C(p,3r+p,3x)=C(p,3r,3x)=3C(p,r,x).
\]
The two identities used here,
\[
 C(v,a+v,b)=C(v,a,b),\qquad
 C(v,3a,3b)=3C(v,a,b),
\]
are unconditional identities proved by cases on $v$. In particular this
argument also avoids conditional parity cases.

\paragraph{Comparison.}
Define
\[
 F(y,a)=I(y,a;u.H(u)),\qquad W(x)=I(x,1;u.u+u),\qquad
 K(x,y)=P(F(y,W(x))).
\]
Three elementary inductions on $y$ give
\[
 H(F(y,a))=F(y,H(a)),\qquad
 F(y,3a)=3F(y,a),\qquad F(y,1)=1.
\]
Their steps use the iterator equation and, respectively, the preceding
induction hypothesis, $H(3u)=3H(u)$, and $H(1)=1$.
Since $W(0)=1$ and $W(Sx)=2W(x)$, we get
\begin{align*}
 K(0,y)&=1,\\
 K(Sx,0)&=0,\\
 K(Sx,Sy)
 &=P(H(F(y,2W(x))))\\
 &=P(F(y,H(2W(x))))\\
 &=P(F(y,3W(x)))
  =P(3F(y,W(x)))=K(x,y).
\end{align*}
Consequently $K(Sx,x)=0$ by induction on $x$.
The intended reading of $K(x,y)$ is the Boolean indicator of $x\leq y$,
but that semantic reading is not needed: only the three displayed equations
and their diagonal consequence are used below.

\paragraph{Saturation and predecessor.}
Put
\[
 B(k,u)=C(K(S(Su),k),u,Su),\qquad
 T(t,k)=I(t,0;u.B(k,u)).
\]
The comparison equations and the unconditional selector identity
$C(v,Sa,Sb)=S(C(v,a,b))$ give
\[
 B(Sk,Su)=S(B(k,u)),\qquad B(S(Sk),0)=1.
\]
Induction on $t$ proves the conjugacy
\begin{equation}
 T(St,S(Sk))=S(T(t,Sk)).\tag{T}\label{eq:saturation-conjugacy}
\end{equation}
At $t=0$ it says $B(S(Sk),0)=1$. For the step, unfold the left side and
rewrite its inner value by the induction hypothesis:
\[
 T(S(St),S(Sk))
 =B(S(Sk),T(St,S(Sk)))
 =B(S(Sk),S(T(t,Sk)))
 =S(T(St,Sk)).
\]
An induction on $n$, applying \eqref{eq:saturation-conjugacy} at $t=k=n$
in its step, now gives $T(n,Sn)=n$.

Finally define the scalar iterator term
\[
 \operatorname{pred}(n)=T(n,n).
\]
Its base equation is immediate, and its successor equation is
\begin{align*}
 \operatorname{pred}(Sn)
 &=B(Sn,T(n,Sn))=B(Sn,n)\\
 &=C(K(S(Sn),Sn),n,Sn)
  =C(K(Sn,n),n,Sn)=C(0,n,Sn)=n.
\end{align*}
Therefore
\[
 \boxed{\operatorname{pred}(0)=0,\qquad
        \operatorname{pred}(Sn)=n.}
\]
Applying this term context to $Sa=Sb$ proves $a=b$.

\paragraph{Dependency conclusion.}
Every induction above has a literal equation as its invariant. Every case
argument proves an unconditional equation and is itself an equation induction
that ignores its induction hypothesis. The construction therefore belongs to
the purely equational iterator fragment. It does not use the full numerical
coverage principle, contextual case splitting, formula induction, successor
injectivity, or a name for an arbitrary term. In a restricted native calculus
with the two iterator equations, equational transport, QE and direct
QI-progress, this entire fragment is available before DI is considered at all.
This is the needed noncircular predecessor lemma for the witness restriction.
